\documentclass{article}
\usepackage{microtype}
\usepackage{graphicx}
\usepackage{subcaption}
\usepackage{booktabs}
\usepackage{xurl}
\usepackage{hyperref}
\usepackage[T1]{fontenc}
\usepackage{icml2026}
\usepackage{amsmath,amssymb,amsfonts,mathtools}
\usepackage{amsthm}
\usepackage{placeins}
\usepackage[nameinlink,capitalize,noabbrev]{cleveref}
\newcommand{\R}{\mathbb{R}}
\newcommand{\Id}{\mathbf{I}}
\newcommand{\norm}[1]{\left\lVert #1\right\rVert}
\newcommand{\Eeq}{E_{\mathrm{eq}}}
\newcommand{\Ehom}{E_{\mathrm{hom,data}}}
\newcommand{\Ehomop}{E_{\mathrm{hom,op}}}
\newcommand{\Erel}{E_{\mathrm{rel,data}}}
\newcommand{\Erelop}{E_{\mathrm{rel,op}}}
\newcommand{\Ecomp}{E_{\mathrm{comp,data}}}
\newcommand{\Ecompop}{E_{\mathrm{comp,op}}}
\newtheorem{proposition}{Proposition}
\newtheorem{corollary}[proposition]{Corollary}
\newtheorem{spectraltheorem}[proposition]{Theorem}
\newtheorem{spectrallemma}[proposition]{Lemma}
\crefname{spectraltheorem}{Theorem}{Theorems}
\crefname{spectrallemma}{Lemma}{Lemmas}
\icmltitlerunning{What Finite-Group Recovery Identifies}
\begin{document}
\twocolumn[
\icmltitle{Beyond Shared Geometry: What Do Recovered Group Actions Identify in Language Models?}
\begin{icmlauthorlist}
\icmlauthor{Anonymous Authors}{anon}
\end{icmlauthorlist}
\icmlaffiliation{anon}{Anonymous Institution}
\icmlcorrespondingauthor{Anonymous Authors}{anonymous@example.com}
\icmlkeywords{representation learning, group actions, equivariance, mechanistic interpretability}
\vskip 0.3in
]
\printAffiliationsAndNotice{}

\begin{abstract}
What do recovered group actions establish about language-model representations? We separate predictive recovery from optimal risk over an exact action class. An analytical example has identical optimal unconstrained regression error and ordinary orbit spectrum, yet different exact-action risks; a spectral bound relates the latter risk to a real-type-constrained invariant-subspace budget. On weekday cycles ($C_7$) and square-corner actions ($D_4$) across six pretrained models, completed controls reveal three empirical findings. Predictive recovery coexists with competitive category and coherently registered alternatives. Predictor rankings change with storage: tested nonlinear predictors do not improve panel means at the linear reference budget, but do at larger budgets. Directly fitting contextual differences improves transport relative to coordinate-fitted operators, while a coherent alternative retains close means. The theory characterizes a whitened exact orthogonal-action risk; the experiments describe separate prediction and transport endpoints on fixed panels. Together they specify what recovery scores distinguish and which competing explanations remain viable.
\end{abstract}

\section{Introduction}
Which additional interpretation is justified when a fitted transformation predicts well, approximately obeys group laws, and outperforms disrupted correspondences? These tests may appear to reinforce one another, yet they address different questions. A useful predictor is not by itself an account of a model's computation; consistency can follow partly from accurate prediction; disrupting a registration does not test whether the specified semantic correspondence is uniquely favored. We study these distinctions through controlled finite actions in contextual language-model representations.

Transformation fitting is established: \citet{lenc2015understanding} study paired maps and network stitching, and \citet{sun2026transformationlaws} develops compatible transformation families, realizability and transport. Equivariance by Contrast identifies jointly learned coordinates and actions under explicit assumptions \citep{schmidt2025equivariance}. Calendar mechanisms are studied by \citet{engels2025notall}; group-operation mechanisms and competing explanations by \citet{nanda2023progress,chughtai2023toy,wu2025unified}. Implementation-family equivalence requires abstraction assumptions \citep{sun2026tracking}. We test which identification claims finite-action controls support; \cref{app:related-work} gives the comparisons.

The protocol fits shared operators to activations of seven weekday states and four corner states across held-out short templates. States are visible in the input, and the language models remain unchanged. Independent fitting removes group-law penalties from the objective, while generator-only fitting withholds action elements. Both succeed predictively. We then compare category predictors, a coherently registered alternative action, compact nonlinear predictors and context-difference fits to determine what these successes distinguish.

\input{generated/main_findings.tex}

Our contribution connects a theoretical distinction with three completed empirical comparisons. The theory shows why regression error and an ordinary spectrum do not determine fixed-dimensional exact-action risk, and bounds that risk using a type-constrained invariant-subspace budget. The empirical findings concern competitive alternative explanations, storage-dependent predictor rankings, and improved context transport under dedicated fitting (\cref{tab:main-findings}). These are complementary results: the new spectral quantity has not been evaluated on the LLM panel. Standard finite-orbit and error-propagation identities explain the interpretation of the controls; their role is supporting analysis rather than additional new theory.

\section{Protocol and Evidence Scope}
\label{sec:setup}
Let $G$ act on states $i$, and let $t$ index contextual templates. For each model, layer, and split, training activations determine mean $\mu$ and orthonormal PCA columns $C\in\R^{d\times q}$. We use
\begin{equation}
 z_{i,t}=(h_{i,t}-\mu)C,\qquad z_{i,t}R_g\approx z_{g(i),t}.
\end{equation}
For row vectors, applying $g_2$ then $g_1$ gives
\begin{equation}
 R_{g_2}R_{g_1}\approx R_{g_1\circ g_2}.
 \label{eq:convention}
\end{equation}
The primary dimension is $q=16$. AFS-D uses matched $q=6$ for seven weekdays and $q=3$ for four corners.

\paragraph{Estimation.}
Ridge fits each element independently, without an intercept, using regularization $\alpha=0.01$. Joint fitting starts from ridge: the data term has coefficient one, and the homomorphism, relation, identity and orthogonality penalties have coefficients $1,1,1,0.05$, respectively. AdamW runs for 1,200 steps, learning rate $0.002$, weight decay $10^{-5}$, and technical seed 42. All transformations use the same fixed training-only PCA. Full objectives appear in \cref{app:methods}.

\paragraph{Evaluation.}
For paired row matrices define
\begin{equation}
 D(A,B)=\frac{1}{m}\sum_{j=1}^m
 \frac{\norm{A_j-B_j}_2^2}{\norm{B_j}_2^2+10^{-12}}.
 \label{eq:D}
\end{equation}
$\Eeq$ averages $D(ZR_g,Z_g)$ over nonidentity actions on held-out templates. $\Ehom$ averages $D(ZR_{g_2}R_{g_1},ZR_{g_1\circ g_2})$ over element pairs. $\Ecomp$ uses the same expression excluding pairs containing the identity; it is a restricted consistency summary, not an independent unseen-element test. AFS-A instead compares generator products with \emph{actual} held-out transformed activations. Relation and operator-space definitions are explicit in \cref{app:metrics}. Data defects are squared relative errors, whereas the reported homomorphism, composition, and relation operator diagnostics use unsquared relative Frobenius norms. Their raw magnitudes cannot be directly compared.

\paragraph{Panel and partitions.}
\label{sec:design}
We use pretrained base Llama-3.2-1B/3B, Llama-3.1-8B, and Qwen2.5-1.5B/3B/7B models at normalized depth slots $0.2,0.4,0.6,0.8$. The $C_7$ action cycles Monday through Sunday. The $D_4$ action rotates and reflects the labels NW, NE, SE, SW. Each domain contains 40 short templates, such as ``The current day is Monday.'' and ``The marked corner is NW.'' The state is visible in the input; no held-out vocabulary, unseen semantic domain, or abstract reasoning task is tested.

For each of five split seeds $(101,303,909,2701,4243)$, 28 templates train PCA and operators, six are diagnostic validation templates, and six are test templates. Complete state orbits stay within their template split. The same 40 templates recur across splits; the seeds represent repeated partitions, not five independently sampled datasets. Activations are mean hidden states over the unique lexical target-token span at transformer block outputs. Exact model revisions, tokenization, layers, prompts, and sensitivity subsets are in \cref{app:methods}.

The evidence comprises a structural recovery panel (FULL and AFS) and subsequent predictor and context-transport comparisons. FULL compares ridge, joint, identity and random controls; AFS tests generator products, disrupted correspondences and rank-matched lexical codes. Later comparisons use separate training-energy normalizations. The appendices map these experiments to their source records, full estimates and preserved stopping decisions.

\paragraph{Statistical scope.}
All empirical comparisons in this paper are descriptive analyses of the recorded fixed panels. Computational replay re-executes inspected designs and supplies no independent confirmation. We retain the recorded means and bootstrap percentile intervals. Historical FULL uses 5,000 draws, base seed 314159 with analysis-specific offsets, and model, depth, split, and template levels, adding random-control seed where relevant. AFS-D keeps depth fixed for its matched contrasts; the separate global-additive experiment resamples split aggregates. The six models are deliberately selected, layers are repeated measurements, and template partitions overlap. These intervals describe resampling sensitivity under the recorded scheme; nominal 95\% does not assert calibrated population coverage. Numerous secondary intervals are not globally multiplicity-adjusted. We make no new significance or equivalence decisions from them; \cref{app:inference} gives their scope.

\section{Prediction Error and Exact-Action Risk}
\label{sec:boundaries}
For a fixed distribution and loss, a predictor's risk upper-bounds the optimum over any function class containing it. It does not determine the optimum over a different class constrained to obey group laws. We make this distinction precise, then relate it to finite-state recovery controls.

\subsection{Same summaries, different exact-action risk}
Consider complete labeled states $s\in S$ in context $c$, with fixed-coordinate observations $z_{c,s}\in\R^{1\times r}$. For a fixed context distribution $P_C$, average uniformly over states and all group elements, including identity and fixed points. For positive-definite second moment $\Sigma=\mathbb E[z^\top z]$, set $x=z\Sigma^{-1/2}$ and $x_g(c,s)=x(c,g\cdot s)$. Let $\eta$ be the minimum normalized squared prediction error for independently chosen matrices $A_g$, and let
\[
 L^\star=\inf_{\rho\in\mathcal R_r(G)}
 \frac1{r|G|}\sum_{g\in G}\mathbb E\norm{x\rho_g-x_g}_2^2,
\]
where $\mathcal R_r(G)$ contains exact row actions with $\rho_g\in O(r)$ and $\rho_g\rho_h=\rho_{hg}$. Both risks use this same whitened metric.

\paragraph{An analytical separation.}
Fix $G=S=C_2\times C_2\times C_7$ acting by addition, uniform states $(a,b,k)$, one context and $r=1$. The observations
\[
 x_A=\frac{(-1)^a+(-1)^b}{\sqrt2},\qquad
 x_B=\sqrt2\cos(2\pi k/7)
\]
have unit second moment, the same optimal regression risk $\eta=1/2$, and the same untyped nonzero orbit spectrum $\{1/2,1/2\}$. Nevertheless,
\[
 \boxed{\ L_A^\star=1,\qquad L_B^\star=2.\ }
\]
The first spectrum occupies two real one-dimensional characters; the second occupies an indivisible real rotation plane. A one-dimensional exact action can retain one of the former but none of the latter. Thus regression error and an ordinary spectrum are insufficient summaries of exact-action recovery. The full labeled observations distinguish these examples; this is a separation of summaries, not nonidentifiability of the full data. \Cref{sp:counterexample} gives the calculation.

\subsection{Bounding the entire exact-action class}
Let $\gamma_r$ be the fraction of observation energy captured by the best group-invariant subspace of dimension at most $r$, respecting real irreducible types. Equivalently, it is the type-constrained spectral budget of the orbit Gram matrix, divided by $r$ (\cref{sp:budget,sp:observable}). \Cref{sp:main} proves
\[
 \boxed{\ 2(1-\gamma_r)\le L^\star\le2(1-\gamma_r^2).\ }
\]
The lower bound holds for every member of $\mathcal R_r(G)$, and some member of the same class attains the upper bound or better. For $\gamma_r<1$, the bounds differ by a factor $1+\gamma_r\le2$; both vanish at $\gamma_r=1$. This connects best invariant approximation with fixed-dimensional prediction risk, using classical averaging and representation tools credited in \cref{sp:positioning}. It is a risk bound, not a demonstrated learning algorithm.

\paragraph{Metric and evidence contract.}
This theoretical risk differs from the historical $D$ and later training-energy-normalized scores. Ridge and finite-step penalized Joint need not be exact orthogonal actions, so their reported scores are not witnesses for this bound. No LLM spectrum or independent-context certificate is reported here. The conservative certificate in \cref{sp:certificate} requires a sampling contract unmet by the present panel. The $r=1$ example establishes a possible distinction, not its empirical magnitude at $r=16$; \cref{sp:type-scale} gives a simple scale restriction for C7/D4.

\subsection{Consequences for finite-state recovery}
Three elementary facts explain the empirical controls; their formal statements and proofs are in \cref{app:theory}. First, any $n$ centered affinely independent state vectors in $n-1$ dimensions admit an exact action of every state permutation (\cref{prop:finite}). Seven generic weekday vectors therefore realize a cycle without group-specific semantic learning. A stable state code plus complementary context features can also share one action across contexts. Held-out templates do not exclude that construction, although arbitrary separate contextual realizations need not share a predictor.

Second, composition residuals depend on prediction errors and operator norms (\cref{prop:propagation}). Independent fitting removes group-law penalties from the objective, but does not show that measured consistency adds identification information beyond accurate prediction. The common-design ridge identity in \cref{app:ridge-identity} makes this dependence explicit.

Third, displacements averaged over a complete uniformly weighted state orbit cancel for every permutation (\cref{prop:additive}). The global additive baseline is therefore an identity predictor by construction. Source-conditioned displacements avoid that cancellation and can compose by telescoping centroid differences (\cref{app:additive-composition}). The comparison below evaluates such alternatives with their extra information stated.

These facts use standard realization and perturbation arguments, related to the broader induced-action questions of \citet{sun2026transformationlaws}. Exact data determine an action only on the supported span; extensions and native implementation remain separate targets. The orientation counterexamples in \cref{app:orientation,app:commutant} further show that equal within-context geometry does not guarantee common transport. They delimit the controls rather than supply a null distribution for the measured LLM errors.

\section{What the Completed Controls Identify}
\label{sec:evidence}
\Cref{tab:main-findings} summarizes the three empirical findings; \cref{tab:identification-summary} retains the detailed endpoint-by-endpoint comparison. Compare losses within each endpoint: historical $D$, compact-coordinate prediction and context-difference transport use different denominators and sampling contracts.

\subsection{Prediction and consistency: two properties, coupled evidence}
Independent ridge attains $\Eeq=0.04830/0.05200$ for $C_7/D_4$, while joint errors are $0.04934/0.05286$. Joint fitting decreases relation defects by $0.01699/0.00728$ at a small direct-fit cost. The maps support approximate predictive action structure without group-law penalties in the independent objective. However, \cref{prop:propagation} explains why this is not automatically an additional semantic discovery beyond accurate prediction. These diagnostics do not establish a comparison matched in prediction error and operator norm.

Generator-only fits predict five excluded nonidentity elements in each group. Their errors are $0.074443/0.056558$, versus directly fitted element errors $0.051124/0.050032$. Correct $D_4$ words outperform matched-length wrong words by $1.097329$ under the recorded contrast. This is meaningful transfer in the operator index; the component states and vocabulary were seen in training. Full structural and generator tables, their historical intervals and depth curves are retained in \cref{app:results}.

\subsection{Invalid registration differs from a coherent alternative}
True-registration homomorphism scores in AFS-B lie below the non-automorphic reference's 2.5\% quantile in all 120 cells per group and metric. Its restricted ensemble excludes the observed registration; no calibrated semantic-null inference is claimed. AFS-C refits Ridge on corrupted training pairs; errors increase against unchanged true held-out targets. Ordinary predictors can degrade similarly, so this establishes sensitivity to supervision rather than semantic specificity.

The completed replay instead consistently supervises an alternative action induced by swapping two state labels. It preserves group laws and fixed-point counts. All-element ridge errors for the specified/coherent correspondence are $0.014897/0.017503$ in $C_7$ and $0.041919/0.041871$ in $D_4$. Each correspondence uses the same underlying source bank and abstract action labels, but its own targets and changed-state evaluation mask. In $D_4$, conjugacy preserves fixed-point counts while potentially changing their identities, so the scored source/action rows need not coincide. This compares the predictability of two registered tasks, rather than two predictors on identical target rows. The single conjugate alternative remains competitive descriptively despite the invalid-registration separation; it does not characterize all semantic alternatives or nonisomorphic representations.

\subsection{Predictor rankings change with information and storage}
\label{sec:characterization}
Rank-matched embedding fits are near exact at $q=6$ for weekdays and $q=3$ for corners; internal errors are higher in the two-model lexical panel. \Cref{prop:finite} explains why this can occur without group-specific semantic learning. Equal rank retains the entire lexical span but can discard internal contextual variation, so it does not match retained information or compare all intrinsic structure at the two locations.

In the $q=16$ replay, the decoded-centroid predictor improves the $C_7$ mean but wins only 9/24 cells; it is worse on the $D_4$ mean and wins in 2/24. It receives training state labels and the exact action table, beyond source-coordinate/action access. This is a strong explanatory control with an information advantage, not a claim of superiority at equal information.

The compact nonlinear comparison fits RFF and Nystr\"om residual predictors using all construction target pairs. Neither improves the group mean at the identity-pruned linear storage budget; Nystr\"om does at $2\times$ storage and RFF at $4\times$ in both groups. \Cref{fig:budget-frontier} shows all five declared nominal budgets; exact means and unit-budget comparisons remain in \cref{tab:reproduction-budgets,tab:reproduction-compact}. Storage includes predictor features, normalizers, coefficients and biases, but excludes shared PCA, model weights and runtime memory. The observed ranking is estimator- and budget-dependent; it is not universal linear superiority or total-compute efficiency. Budgets were not selected retrospectively as deployment operating points.

\begin{figure*}[t]
\centering
\includegraphics[width=\textwidth]{figures/budget_frontier.pdf}
\caption{\textbf{Predictor rankings change with storage.} Recorded direct-prediction means for all nonidentity actions on changed states under the specified semantic registration, over 24 fixed model/readout cells per domain. Losses use the replay training-coordinate-energy denominator; lower is better. The horizontal axis is the nominal incremental parameter-storage budget relative to identity-pruned ridge, not runtime memory or statistical capacity. All five declared budgets and all three budgeted predictor families are shown. Nystr\"om first improves the group mean over ridge at $2\times$, and RFF at $4\times$, in both groups. Lines connect recorded aggregates on logarithmic axes; they are not fitted scaling laws or confidence bands. This direct-prediction endpoint is distinct from generator-only composition.}
\label{fig:budget-frontier}
\end{figure*}

\subsection{Fitting context differences improves transport}
Fitting context differences directly yields losses $0.03426/0.11639$ for $C_7/D_4$, versus $0.04714/0.13230$ for applying coordinate-fitted operators to those differences. The coherent action gives $0.03446/0.11717$ under the same direct fitting procedure and registration-specific scoring just described. The improvement therefore identifies a useful fitting objective, while close means for distinct registered tasks do not establish semantic preference or equivalence. On jointly new $S_3$ contexts/entities, direct transport loss is $0.016290$ versus identity $0.016444$, a small absolute improvement below the inherited $0.01$ practical margin. Complete normalization and information conditions are in \cref{app:characterization,app:reproduction}.

\paragraph{Behavioral scope.}
Historical activation edits supply limited behavioral evidence: $C_7$ target logit improves over identity by 1.129, but the mean effect is smaller than that of the source-conditioned additive control. Correct $D_4$ ordering has lower final-logit distance than the matched wrong ordering. These are token/logit outcomes, not a certified native algorithm. The additive control uses full-hidden displacements and source identity, whereas the operator edits a PCA subspace; the historical mean-overwrite hook lacks a matched zero-delta sham. The separate native program branch failed clean competence and its confirmation remains blocked. Full endpoints and limits, including the unchanged failure records, are in \cref{app:behavior,app:reproduction}.

\section{Implications for Evaluating Recovered Actions}
\label{sec:reading-protocol}
The findings motivate four complementary questions, conditional on states, contexts, coordinates, estimator and outcome. They organize interpretation rather than define a calibrated certification test or four statistically independent sources of evidence.

\paragraph{1. Predictive recovery.}
State what was held out: contexts, action elements, state identities or domains. Disclose comparator supervision and storage. Our generator result demonstrates why the distinction matters: transfer to an excluded action element can still use a completely familiar state vocabulary.

\paragraph{2. Action consistency.}
Specify the group-law residual, its metric and whether fitting encouraged it. The propagation identity relates consistency to prediction errors and operator norms. To claim additional identification information, compare predictors under conditions that account for those quantities; independent objectives alone do not establish independent evidence.

\paragraph{3. Preferential semantic alignment.}
Compare coherent alternatives as well as invalid registrations, making information, capacity, scoring rows, target definitions and predictive difficulty explicit. Our comparisons show why these controls answer different questions. Preference over selected alternatives would narrow the interpretation, rather than identify a unique semantic mechanism.

\paragraph{4. Native causal use.}
For a claim about successful native computation, establish clean-task competence and intervention fidelity, then distinguish sufficiency from necessity for the proposed mechanism. The historical edits illustrate the gap between moving an output score and identifying the computation used by the unedited model.

For a new application, specify the intended claim and competing explanation before choosing a diagnostic. The useful outcome is the interpretation changed by the comparison, including unfavorable results, rather than the number of positive scores obtained.

\section{Discussion and Remaining Scope}
The experiments show that predictable finite actions coexist with competitive explanatory alternatives. Storage changes predictor rankings, and fitting context differences directly improves transport without establishing preference for the specified semantic registration. The analytical identities and control definitions limit what additional identification information is justified by combining prediction, group-law and disrupted-control scores. The analytical separation shows that identical regression risk and ordinary spectra can conceal different exact-action risks. The spectral bound characterizes that class optimum under its own metric; connecting it empirically to the LLM panel remains open. These results specify what the tested controls change in an interpretation, rather than introduce a new representation theory or supersede competent-task causal studies.

The main open issue is robustness beyond the present construction. C7/D4 use small seen-state vocabularies; repeated partitions reuse templates, and the later study uses a fixed grammar generator. S3 includes new entities and contexts but yields only small gains over identity. Models and layers are selected panels, not independent population samples. Structural power remained limited and native confirmation is blocked. The appendix retains secondary peer, cross-model, product and trained-interface results with their distinct endpoints and original failure records.

A bounded next step is to check whether saved, complete, same-coordinate banks permit the finite-bank spectral diagnostic in \cref{sp:observable}. Such a retrospective analysis could connect the theoretical quantities to this panel, but would not provide new-context validation. An independently authored, prospectively frozen context bank would instead test whether the principal comparisons survive a different context construction. A focused follow-up could compare specified and multiple coherent registrations on one common complete source/action grid, with fixed information, fitting rules and normalization. That would be a new evaluation contract, unlike the registration-specific masks used here. It would strengthen evidence about generalization and registration preference, not by itself exclude state coding or establish native causal use. It has not been run and is not needed to state the present fixed-panel findings.

\section{Conclusion}
Across the declared language-model panel, finite-action predictors transfer to held-out contexts and generator-composed elements. Their interpretation changes when category information, predictor storage and context-difference fitting are made explicit. The analytical example and risk bound additionally distinguish unconstrained prediction from the best exact orthogonal action of a fixed dimension. In the empirical panel, dedicated fitting improves transport while the tested coherent alternative remains descriptively competitive. This combination of predictive success and viable alternatives motivates evaluating recovery by the explanation each control distinguishes, with semantic preference and native causal use treated as separate empirical questions.

\FloatBarrier
\label{sec:main-end}

\section*{Impact Statement}
This work develops methods for evaluating structure in language-model representations. Better diagnostics may aid interpretability, while overinterpreting probes could produce misleading claims about model capabilities or mechanisms. The reported controlled tasks do not evaluate demographic fairness or establish safety guarantees. We identify no additional specific societal consequences requiring separate discussion.

\bibliography{references}
\bibliographystyle{icml2026}
\appendix
\onecolumn
\section{Proofs and the Identification Problem}
\label{app:theory}
\label{sec:appendix}
\subsection{Supporting finite-state statements}
The following elementary statements support Section~\ref{sec:boundaries}; they are not presented as new general theory. Their original mathematical content is retained.

\begin{proposition}[Finite-orbit realizability]
\label{prop:finite}
Let $X\in\R^{n\times(n-1)}$ have centered rows and rank $n-1$. For a permutation $g$, define $(P_gX)_i=X_{g(i)}$. There is a unique invertible $R_g=X^\dagger P_gX$ satisfying $XR_g=P_gX$, and $R_gR_h=R_{h\circ g}$.
\end{proposition}

Write $e_g(i,t)=z_{i,t}R_g-z_{g(i),t}$.
\begin{proposition}[Prediction-error propagation]
\label{prop:propagation}
For $p=g_1\circ g_2$, the composition residual is
\begin{align}
 z_{i,t}(R_{g_2}R_{g_1}-R_p)
 &=e_{g_2}(i,t)R_{g_1} \nonumber\\
 &\quad+e_{g_1}(g_2(i),t)-e_p(i,t).
\end{align}
Its norm is at most $\norm{e_{g_2}}_2\norm{R_{g_1}}_{\rm op}+\norm{e_{g_1}(g_2(i),t)}_2+\norm{e_p}_2$.
\end{proposition}

\begin{proposition}[Global-field cancellation]
\label{prop:additive}
For any permutation $g$ of a complete uniformly weighted orbit, $\sum_i(z_{g(i),t}-z_{i,t})=0$. Therefore the train-fitted global displacement $v_g=\mathbb{E}_{i,t}[z_{g(i),t}-z_{i,t}]$ is zero for every element of both tested groups, up to numerical precision.
\end{proposition}

\subsection{Finite-orbit realizability}
\begin{proof}[Proof of \cref{prop:finite}]
Centering implies $\operatorname{col}(X)\subseteq\mathbf{1}^{\perp}$. Both spaces have dimension $n-1$, so they are equal. Every permutation matrix preserves $\mathbf{1}^{\perp}$; hence $P_gX$ lies in the column space of $X$. Full column rank gives the unique solution $R_g=(X^\top X)^{-1}X^\top P_gX$. Our definition implies $P_gP_h=P_{h\circ g}$, and consequently
\[
 XR_gR_h=P_gXR_h=P_gP_hX=P_{h\circ g}X=XR_{h\circ g}.
\]
Uniqueness establishes the multiplication rule. Taking $h=g^{-1}$ proves $R_gR_{g^{-1}}=R_e=\Id$ and the reverse product similarly, so every $R_g$ is invertible.
\end{proof}
If $n$ lexical vectors are affinely independent, their centered span has dimension $n-1$. Keeping all these PCA dimensions makes only an invertible coordinate change, so the proposition still applies. Neither regular polygon geometry nor pretrained knowledge of a named group is required. The induced matrices need not be orthogonal in the measured Euclidean metric. Their exact abstract action and their geometric shape are separate questions.

This interpolation issue also limits representation comparison: invariance to arbitrary invertible changes can make sufficiently high-dimensional sample representations indistinguishable \citep{kornblith2019similarity}. That observation concerns comparison statistics; it does not establish out-of-context transformation prediction.

\subsection{Separate contextual realizations need not be shared}
Consider centered three-state coordinates
\[
 X=\begin{bmatrix}1&0\\0&1\\-1&-1\end{bmatrix},\qquad
 R=\begin{bmatrix}0&1\\-1&-1\end{bmatrix}.
\]
For the cycle $g:1\mapsto2\mapsto3\mapsto1$, $XR=P_gX$. A second contextual copy $Y=XA$, with $A=\operatorname{diag}(2,1)$, has forced operator
\[
 R'=A^{-1}RA=\begin{bmatrix}0&1/2\\-2&-1\end{bmatrix}\ne R.
\]
A shared operator $S$ would satisfy both $XS=P_gX$ and $YS=P_gY$. Full column rank forces $S=R$ from the first and $S=R'$ from the second, a contradiction. This gives a precise reason why training one operator family and evaluating held-out contexts is stronger than recovering a separate action on each finite orbit.

\subsection{Shared contextual realizations can still reflect state identity}
The preceding counterexample is a non-guarantee, not evidence that context uniformity uniquely identifies semantic computation. Let $x_i$ be any centered full-affine-rank state code, and let $c_t$ be any context vector. In the direct-sum representation
\[
 z_{i,t}=(x_i,c_t),\qquad
 R_g=\begin{bmatrix}A_g&0\\0&\Id\end{bmatrix},
\]
where $x_iA_g=x_{g(i)}$, one shared family maps all contexts exactly and satisfies all group relations. No contextual calculation of a weekday shift or geometric rotation is necessary for this construction. The same state code admits consistently registered alternative permutation actions. A global invertible coordinate change preserves this example up to conjugacy, so visible separation of state and context coordinates is unnecessary.

This construction is a sufficiency example, not a claim that the tested activations have this exact form. Historical AFS-B changes operator labels inconsistently with the registered multiplication table, while AFS-C changes correspondence across training contexts; both can damage a stable state code. The later characterization and replay explicitly evaluate category predictors and a coherent alternative registration (\cref{sec:characterization,app:reproduction}). Their competitive results leave stable state coding compatible with the observations. These comparisons do not establish semantic specificity or constitute a conditional test jointly matched in prediction error, operator norm, information and capacity.

\subsection{Shared contextual actions as a subspace test}
\label{app:contextual-realization}
Stack contextual rows into $Z\in\R^{N\times q}$ and let $P_g$ permute states within each complete template. At any rank, standard least-squares projection gives
\[
 \min_R\norm{ZR-P_gZ}_F^2
 =\norm{(\Id-ZZ^\dagger)P_gZ}_F^2.
\]
Indeed, $ZZ^\dagger P_gZ$ is the orthogonal projection of the target matrix onto $\operatorname{col}(Z)$, leaving the displayed orthogonal residual. Thus an exact shared map exists if and only if this residual vanishes, equivalently $P_g\operatorname{col}(Z)\subseteq\operatorname{col}(Z)$. When $Z$ has full column rank, the exact solution $R_g=Z^\dagger P_gZ$ is unique; group-wide existence then gives the opposite multiplication convention of \cref{prop:finite}.

Equivalently, let $K=\ker(Z^\top)$, the space of linear dependencies among rows. Exact realizability is equivalent to
\[
 P_g^\top K\subseteq K.
\]
Necessity follows from $a^\top Z=0\Rightarrow a^\top P_gZ=a^\top ZR_g=0$. Conversely, this condition makes the rule $a^\top Z\mapsto a^\top P_gZ$ well defined on the row span; it therefore defines a linear map there. For a whole group the induced maps are invertible and obey the row-action law. This is invariance of the \emph{left} kernel; preservation of only equal-row pairs, or of the right kernel $\ker Z$, is not this linear-realizability criterion. For example, the three distinct centered scalar states $(-1,0,1)$ have no exact linear realization of their three-cycle, since $(0,1,-1)$ is not a scalar multiple of $(-1,0,1)$.

When $Z$ is rank deficient, only the action on its row span is identified. The minimum-norm solution has $R_e=Z^\dagger Z$, a projector rather than the ambient identity, and need not be invertible. If all the subspace conditions hold, a coherent invertible extension is
\[
 \widetilde R_g=Z^\dagger P_gZ+(\Id-Z^\dagger Z),
\]
which acts as the induced action on the row span and as identity on its orthogonal complement. It satisfies $\widetilde R_e=\Id$ and $\widetilde R_g\widetilde R_h=\widetilde R_{h\circ g}$. Its unobserved complementary action is a choice, not a property identified by the data. This uniqueness concerns fixed coordinates on the finite observed bank: different ambient matrices can agree on its span, while new contexts can impose additional constraints.

This elementary criterion connects shared contextual fitting with subspace preservation. It is an unregularized, unnormalized characterization, not a formula for the recorded row-normalized test error or evidence of population generalization. Ridge instead returns
\[
 R_g^{(\alpha)}=(Z^\top Z+\alpha\Id)^{-1}Z^\top P_gZ.
\]
If an exact $R_g^\star$ exists, this is $B_\alpha R_g^\star$, where $B_\alpha=(Z^\top Z+\alpha\Id)^{-1}Z^\top Z$. In particular, the independently fitted identity is $B_\alpha$, not necessarily $\Id$. Shrinkage can therefore introduce power and relation errors even in an exactly realizable system. Joint relation/identity penalties target precisely such deviations. The current comparison does not separate those regularization effects from any additional semantic information.

\subsection{An exact common-design ridge identity}
\label{app:ridge-identity}
Let every element use the same source matrix $Z\in\R^{N\times q}$, regularization $\alpha>0$, and permutation targets $P_gZ$. Write
\begin{align*}
 L_\alpha&=(Z^\top Z+\alpha\Id)^{-1}Z^\top,\\
 K_\alpha&=\Id-ZL_\alpha,\qquad R_g=L_\alpha P_gZ.
\end{align*}
Then the paper's convention $P_gP_h=P_{h\circ g}$ gives the exact identities
\begin{align}
 ZR_g-P_gZ&=-K_\alpha P_gZ,\label{eq:ridge-fit-identity}\\
 R_gR_h-R_{h\circ g}&=-L_\alpha P_gK_\alpha P_hZ.
 \label{eq:ridge-law-identity}
\end{align}
To prove them, substitute $ZL_\alpha=\Id-K_\alpha$ into $ZR_g$ and into $L_\alpha P_g(ZL_\alpha)P_hZ$, respectively. No exact realizability or full-rank assumption is needed. For positive $\alpha$, $K_\alpha$ is a ridge residual operator, not in general an orthogonal projector. These equalities make the dependence between fit and composition defects explicit for a common design. They are not independent significance tests, and they do not apply unchanged to element-specific training subsets, weights, regularization, or independently centered targets. Their unnormalized matrices must not be identified with the row-normalized test metric in \cref{eq:D}.

\subsection{Equal within-context geometry need not permit shared transport}
\label{app:orientation}
We give a restricted analytical construction, using standard orthogonal representation and averaging tools \citep{serre1977linear}. It specifies an identification boundary; it is not an assumed model of the measured activations or an additional calibration experiment.

Let $V\in\R^{n\times r}$ have centered rows, $V^\top V/n=\Id_r$, and $VQ_g=P_gV$, where $Q_g\in O(r)$ and $Q_gQ_h=Q_{h\circ g}$. For $d\geq r$, choose $A_c\in\R^{d\times r}$ with $A_c^\top A_c=\Id_r$ and set $X_c=VA_c^\top$. State count $n$, orbit rank $r$, and ambient dimension $d$ are distinct. Every context has identical within-context Gram matrix $X_cX_c^\top=VV^\top$ and an exact local orthogonal action
\[
 R_{c,g}=A_cQ_gA_c^\top+(\Id_d-A_cA_c^\top).
\]
Indeed $X_cR_{c,g}=P_gX_c$, and the two orthogonal subspaces carry $Q_g$ and the identity. Cross-context inner products are not fixed by this construction.

\begin{proposition}[Independent-orientation risk]
\label{prop:orientation}
For a nonempty $S\subseteq G$, define
\[
 \ell_c(R)=\frac{1}{nr|S|}\sum_{g\in S}\norm{X_cR_g-P_gX_c}_F^2.
\]
If a fresh $A_c$ is Haar-uniform on the Stiefel manifold and independent of the fitted family $\widehat R$, then, conditional on that fit,
\begin{equation}
\begin{aligned}
 &\mathbb E[\ell_c(\widehat R)\mid\widehat R]\\
 &\quad=\beta_S+\frac{1}{d|S|}\sum_{g\in S}
 \norm{\widehat R_g-\frac{\chi(g)}r\Id_d}_F^2,
\end{aligned}
\end{equation}
where $\chi(g)=\operatorname{tr}(Q_g)$ and
\[
 \beta_S=1-\frac1{|S|}\sum_{g\in S}\left(\frac{\chi(g)}r\right)^2.
\]
Thus $\beta_S$ is the minimum expected risk among unrestricted shared matrix predictors in this model.
\end{proposition}
\begin{proof}
Whitening gives $\norm{X_cR-P_gX_c}_F^2/(nr)=\norm{A_c^\top R-Q_gA_c^\top}_F^2/r$. Haar invariance and traces imply $\mathbb E[AA^\top]=(r/d)\Id_d$ and $\mathbb E[AQ_gA^\top]=(\chi(g)/d)\Id_d$. Expanding the square yields
\begin{align*}
 &\frac1d\norm R_F^2+1-\frac{2\chi(g)}{rd}\operatorname{tr}R\\
 &\quad=1-\frac{\chi(g)^2}{r^2}
   +\frac1d\norm{R-\frac{\chi(g)}r\Id_d}_F^2.
\end{align*}
Average over $S$. Independence allows conditioning on the learned matrices; the scalar family attains the bound, although it need not itself be a group representation.
\end{proof}
For the centered seven-state cycle, $r=6$, $\chi(e)=6$, and $\chi(g)=-1$ for every other element. Hence $\beta_G=5/6$ and $\beta_{G\setminus\{e\}}=35/36$. For the centered four-corner $D_4$ action, $r=3$ and each nonidentity character is $+1$ or $-1$, giving $7/9$ and $8/9$, respectively. These constants apply only to the specified whitened-orbit risk and independence hypothesis. Comparing them numerically with the reported $\Eeq$ would mix different distributions and normalizations. No Haar randomization test or empirical rejection of this model is claimed.

\subsection{A commutant criterion in the full-rank construction}
\label{app:commutant}
Specialize the preceding construction to $d=r$, writing $A_c=O_c\in O(r)$. Orientations can now have any distribution. Define
\begin{align*}
 \mathcal C_Q&=\{B:BQ_g=Q_gB\ \forall g\},\\
 \Pi_Q(B)&=\frac1{|G|}\sum_g Q_gBQ_g^\top.
\end{align*}
The group average is the Frobenius orthogonal projection onto $\mathcal C_Q$: reindexing the group shows its image commutes with every $Q_g$, it fixes every commuting matrix, and inverse closure makes the average self-adjoint.

\begin{proposition}[Relative-orientation obstruction]
\label{prop:commutant}
For independent copies $c,c'$ of the context, $U=O_c^\top O_{c'}$, and $S=G$,
\[
 \inf_{\{R_g\}}\mathbb E_c\ell_c(R)
 =\frac1r\mathbb E_{c,c'}\norm{U-\Pi_Q(U)}_F^2.
\]
Consequently an exact shared action exists almost surely precisely when the relative orientations commute with every $Q_g$ almost surely.
\end{proposition}
\begin{proof}
The local action is $R_{c,g}=O_cQ_gO_c^\top$, and the loss for element $g$ is $\norm{R_g-R_{c,g}}_F^2/r$. The unrestricted minimizer is $\mathbb E_cR_{c,g}$. The pairwise variance identity therefore gives optimum risk
\[
 \frac1{2r|G|}\mathbb E_{c,c'}\sum_g\norm{Q_gU-UQ_g}_F^2.
\]
Expansion and the projection identity give $|G|^{-1}\sum_g\norm{Q_gU-UQ_g}_F^2=2\norm{U-\Pi_Q(U)}_F^2$. Nonnegativity proves the equivalence; when the risk is zero the common matrices equal a local action and inherit its group law. At positive risk, the mean matrices need not form a representation.
\end{proof}
Large context rotations within the commutant can therefore leave the action unchanged. This is a full-rank, whitened, noiseless characterization; it is not a characterization of every contextual representation. For $m$ iid contexts in this same construction, define the unregularized mean of the local matrices and its empirical risk by
\[
 \widehat R_{m,g}=\frac1m\sum_{j=1}^mR_{c_j,g},\qquad
 \widehat L_m=\frac1m\sum_{j=1}^m\ell_{c_j}(\widehat R_m),\qquad
 b=\inf_{\{R_g\}}\mathbb E_c\ell_c(R).
\]
For an independent $c_{\rm new}$, the variance identity gives
\[
 \mathbb E\widehat L_m=(1-1/m)b,\qquad
 \mathbb E\ell_{c_{\rm new}}(\widehat R_m)=(1+1/m)b.
\]
Both expectations include the sampled training contexts; the second also averages over the fresh context. These are sample-mean identities, not a learning-curve forecast for ridge fits to the actual $q=16$ activations. The proofs supply analytical checks on interpretation, not a claim of new general representation theory or a validation of the construction in the language models.

\subsection{Prediction-error propagation}
\begin{proof}[Proof of \cref{prop:propagation}]
Insert and subtract $z_{g_2(i),t}R_{g_1}$ and $z_{g_1(g_2(i)),t}$. Since $p(i)=g_1(g_2(i))$,
\begin{align*}
 z_{i,t}R_{g_2}R_{g_1}-z_{i,t}R_p
 &=\bigl(z_{i,t}R_{g_2}-z_{g_2(i),t}\bigr)R_{g_1}\\
 &\quad+\bigl(z_{g_2(i),t}R_{g_1}-z_{p(i),t}\bigr)
       -\bigl(z_{i,t}R_p-z_{p(i),t}\bigr).
\end{align*}
This is the stated identity. The triangle inequality and $\norm{vA}_2\leq\norm{v}_2\norm{A}_{\rm op}$ give the norm bound.
\end{proof}
For example, if every unnormalized prediction error is at most $\delta$ on a group-closed evaluation set and $\norm{R_g}_{\rm op}\leq M$, every pairwise residual is at most $(M+2)\delta$. A corresponding squared bound follows from $(a+b+c)^2\leq3(a^2+b^2+c^2)$. If additionally every comparison target satisfies $\norm{z_{i,t}R_p}_2^2+10^{-12}\geq b>0$, each corresponding normalized defect is bounded by $(M+2)^2\delta^2/b$. Neither a useful $b$ nor a uniform $\delta$ follows merely from a small average direct error. Thus this is not a numerical bound on the reported $D$ without additionally accounting for its sample-specific denominators. Nor does it certify small composition defects when norms are uncontrolled. It establishes a dependence that must be addressed before claiming algebraic consistency beyond what prediction already entails. We have not performed an error-matched conditional null analysis in the present experiments.

\subsection{Global additive cancellation}
\begin{proof}[Proof of \cref{prop:additive}]
Since $g$ is a bijection, $\{g(i):i=1,\ldots,n\}=\{1,\ldots,n\}$. Thus
\[
 \sum_{i=1}^n z_{g(i),t}=\sum_{i=1}^n z_{i,t}.
\]
Subtracting gives zero separately in every template. Averaging over any training templates preserves zero. Centering and PCA are unnecessary for this cancellation.
\end{proof}
The conclusion holds for every element of both tested groups, including reflections, and for any representation dimension. More generally it holds for a weighting invariant under the action. An incomplete or non-invariantly weighted orbit need not cancel. Those designs were not the tested estimator. A vector conditioned on source state, $v_{g,i}=\mathbb{E}_t[h_{g(i),t}-h_{i,t}]$, also need not cancel because the source average is removed.

\subsection{State-conditioned additive actions can compose}
\label{app:additive-composition}
For training-only state means $\mu_i=\mathbb{E}_{t\in\mathcal T_{\rm train}}h_{i,t}$ and the complete balanced training orbits, the existing comparator satisfies
\[
 v_{g,i}=\mathbb{E}_t(h_{g(i),t}-h_{i,t})=\mu_{g(i)}-\mu_i.
\]
To compose these edits, retain the source-state index and update it after each action. The augmented-state map is
\[
 A_g(h,i)=(h+\mu_{g(i)}-\mu_i,\ g(i)).
\]
For $p=g_1\circ g_2$,
\begin{align*}
 (A_{g_1}\circ A_{g_2})(h,i)
 &=\bigl(h+(\mu_{g_2(i)}-\mu_i)
             +(\mu_{p(i)}-\mu_{g_2(i)}),p(i)\bigr)\\
 &=\bigl(h+\mu_{p(i)}-\mu_i,p(i)\bigr)=A_p(h,i).
\end{align*}
The means telescope. For noncommuting actions the two orders may have different target means, so this source-conditioned additive rule can preserve order as well as exact composition. Its composition law is on the augmented state carrying the source identity; it is not a constant vector update or an identified native neural mechanism. The behavioral experiment supplies this identity as an oracle condition and fits vectors in the full hidden dimension. The historical study does not provide a structural comparison in the same PCA subspace with matched identity information. Subsequent excluded development comparisons informed the separate characterization study. Its returned coordinate-level results have since been audited; these label-informed controls retain their declared information conditions and do not identify a native mechanism. The identity also shows why global additive cancellation does not establish an intrinsic advantage of multiplicative notation: every linear map obeys $zR_g=z+z(R_g-\Id)$.

The preceding exact composition guarantee assumes that the state index is carried forward. A stateless rule that re-decodes every intermediate vector need not compose: with centroids $\mu_0=0$, $\mu_1=2$ and the swap $s$, nearest-centroid decoding gives $F_s(-2)=0$ and $F_s(0)=2$, so $F_s(F_s(-2))\ne-2$ although $s^2=e$. Perfect decoding on clean evaluation inputs alone does not certify decoding after edits. Likewise, looking up the known composite action after a single input decode is not an empirical test of repeated learned primitive edits.

\subsection{Faithful actions need not be free}
\label{app:faithful-not-free}
Faithfulness means that no nonidentity group element fixes every state. Freeness means that no nonidentity element fixes even one state. The four-corner $D_4$ action is faithful but not free: each corner has a stabilizer of size two, so a single source corner has only four distinguishable endpoints for eight group elements. Aggregating a full orbit can distinguish the transformations; an arbitrary single endpoint cannot. Under $r(i)=i+1$ and $s(i)=1-i$ modulo four, $srs(i)=i-1$ and $rsr(i)=1-i$ coincide at $i=1,3$, whereas $rs(i)=2-i$ and $sr(i)=-i$ differ at all four corners. Order comparisons must therefore report their endpoint-discriminating state set. Success on this action is not evidence of a free eight-state action or of native group multiplication.

\subsection{Coordinate equivalence and the metric of approximation}
\label{app:coordinates}
For an invertible coordinate change $S$, set $z'=zS$ and $R'_g=S^{-1}R_gS$. Then
\[
 z'R'_g-z'_{g}= (zR_g-z_g)S,\qquad
 R'_{g_2}R'_{g_1}-R'_{p}
 =S^{-1}(R_{g_2}R_{g_1}-R_p)S.
\]
Consequently, exact realization and exact group identities are invariant under simultaneous conjugacy. This does not make an arbitrary approximate score coordinate invariant. Orthogonal $S$ preserves the Euclidean row norms and Frobenius norms used here, and therefore the corresponding recorded defects. General invertible $S$ can change their values, their conditioning, and the weighting of representation directions. Euclidean ridge penalties and PCA truncation likewise encode a chosen geometry; refitting after a nonorthogonal change need not yield the conjugated original estimate. Changing the retained PCA subspace is not merely rotating its coordinates.

The semantic action group acting on input states is distinct from the coordinate-equivalence group acting on the representation description \citep{yao2026groupaction}. Our coordinates and semantic labels are fixed. Context-specific realignments would change the shared-operator question; coherent state relabelings define alternative registered tasks, rather than a free equivalence when testing semantic preference. The present numerical claims concern the specified unwhitened PCA metric. They are not invariant quantitative claims under arbitrary reparameterizations. For an unnormalized activation-conditioned matrix defect $D$, the second-moment identity
\[
 \mathbb{E}\norm{zD}_2^2=\operatorname{tr}(D^\top\Sigma D),\qquad
 \Sigma=\mathbb{E}[z^\top z],
\]
shows how covariance weights directions. A small data-space value can reflect low variance rather than an exactly unsupported subspace. Our actual $D(A,B)$ additionally normalizes each row by its target norm, and the operator diagnostics use different powers and, in the C7 relation case, different aggregands. Their numerical comparison cannot by itself establish a sharp supported/unsupported dichotomy.


For an exact finite group action, an invariant positive-definite metric always exists: $H=|G|^{-1}\sum_gR_gR_g^\top$ satisfies $R_hHR_h^\top=H$. Choosing coordinates with $SS^\top=H$ makes the conjugated matrices orthogonal. This fact does not make the measured Euclidean metric invariant, and approximate fits need not admit this exact construction. It further separates abstract algebraic realizability from a preferred geometry of the observed activations.

The same distinction applies to optional character diagnostics in the released code. The centered full-rank $n$-state representation has character $\chi(g)=\#\{i:g(i)=i\}-1$, whereas the full permutation representation has character equal to the number of fixed states. Comparing normalized traces against the latter reference does not measure distance to every valid representation, particularly after centering or changing dimension. Character-inner-product expressions computed from approximately consistent fitted matrices are descriptive real numbers; they are not certified integer irreducible multiplicities.

\section{Fixed-Dimensional Risk Bounds for Exact Orthogonal Actions}
\label{app:spectral-risk}

This appendix separates the best risk in an entire action class from the error of selected fitted operators. The class and metric below are narrower than arbitrary linear prediction: they impose exact group laws and orthogonality after population second-moment whitening. The result is a mathematical extension, not an additional empirical result or a certificate for the recorded FULL/AFS scores.

\subsection{Observation contract and invariant-subspace budget}
\label{sp:setup}
Let a finite group $G$, of order $m$, act on a finite state set $S$, of size $p$; the action may be nonfree or nonfaithful. A context $c$ supplies the complete labeled block $\{z_{c,s}:s\in S\}$ in fixed coordinates $z_{c,s}\in\R^{1\times r}$, with $r\ge1$. Expectations use $C\sim P_C$ and an independent uniform $s\in S$; group averages include every element, including the identity. Assume $\Sigma=\mathbb E[z_{C,s}^{\top}z_{C,s}]\succ0$, and put $x=z\Sigma^{-1/2}$, so $\mathbb E[x^{\top}x]=\Id_r$. This is a second moment; zero mean is not required. The group product $gh$ means first $h$, then $g$, giving the row-action convention
\[
 \mathcal R_r(G)=\{\rho:G\to O(r):\rho_e=\Id_r,\ \rho_g\rho_h=\rho_{hg}\}.
\]
Define two population risks, with $x_g(c,s)=x(c,g\cdot s)$:
\begin{equation}
 \eta=\inf_{\{A_g\}}\frac1{rm}\sum_g\mathbb E\norm{xA_g-x_g}_2^2,
 \qquad L^\star=\inf_{\rho\in\mathcal R_r(G)}L(\rho),\quad
 L(\rho)=\frac1{rm}\sum_g\mathbb E\norm{x\rho_g-x_g}_2^2.
 \label{sp:risks}
\end{equation}
The first class allows independent unconstrained $A_g\in\R^{r\times r}$; the second requires one exact $r$-dimensional orthogonal action shared across contexts. In raw coordinates $R_g=\Sigma^{-1/2}\rho_g\Sigma^{1/2}$, its loss is
\[
 \frac1{rm}\sum_g\mathbb E\norm{(zR_g-z_g)\Sigma^{-1/2}}_2^2.
\]
This is neither the row-normalized metric $D$ in \cref{eq:D} nor the later training-energy normalization. Exact zero-error actions preserve $\Sigma$ and become orthogonal after whitening, but positive-error optima over all $GL(r)$ need not do so. Accordingly, a lower bound below excludes only the class in \cref{sp:risks}. Independently fitted Ridge and finite-step penalized Joint generally do not belong to that exact class. Their recorded errors are neither covered by its lower bound nor upper-bound witnesses for $L^\star$.

On the real Hilbert space $\mathcal H=L^2(P_C\times\mathrm{Unif}(S))$, define
\[
 (Ja)(c,s)=x_{c,s}a,\qquad (U_gf)(c,s)=f(c,g\cdot s),\qquad
 P=JJ^*,\qquad T=\frac1m\sum_gU_gPU_g^*.
\]
Then $J^*J=\Id_r$, $P$ is a rank-$r$ orthogonal projection, and $U_gU_h=U_{hg}$ with each $U_g$ orthogonal. The group average $T$ commutes with every $U_g$, satisfies $0\preceq T\preceq\Id$ and $\operatorname{tr}T=r$, and is supported on
\[
 \mathcal W=\operatorname{span}\{U_gJa:g\in G,\ a\in\R^r\},\qquad \dim\mathcal W\le mr.
\]
Define
\begin{equation}
 C_r(T)=\max_{Q}\operatorname{tr}(TQ),\qquad \gamma_r=C_r(T)/r,
 \quad
 \begin{gathered}
 Q^2=Q=Q^*,\quad [Q,U_g]=0\ \text{for all }g,\\[-2pt]
 \operatorname{rank}Q\le r,\quad\operatorname{range}Q\subseteq\mathcal W.
 \end{gathered}
 \label{sp:budget}
\end{equation}
Thus $0\le\gamma_r\le1$. Since invariant $Q$ satisfies $\operatorname{tr}(PQ)=\operatorname{tr}(TQ)$,
\[
 r-C_r(T)=\min_Q\norm{(\Id-Q)J}_{\mathrm{HS}}^2.
\]
This is a best invariant-subspace approximation quantity, not a new notion of identifiability. Spectral constructions for invariant approximation have precedents \citep{barbieri2020approximation}; the connection to the exact-action risk is established below.

The rank constraint must respect \emph{real} irreducible types. If type $\tau$ has real dimension $d_\tau$ and multiplicity $m_\tau$, the corresponding self-adjoint multiplicity operator has eigenvalues $\lambda_{\tau,1}\ge\cdots\ge\lambda_{\tau,m_\tau}\ge0$, each repeated $d_\tau$ times in the real space. With integer multiplicities $k_\tau$, standard real representation theory \citep{serre1977linear} gives
\begin{equation}
 C_r(T)=\max_{\substack{0\le k_\tau\le m_\tau\\\sum_\tau d_\tau k_\tau\le r}}
 \sum_\tau d_\tau\sum_{j=1}^{k_\tau}\lambda_{\tau,j}.
 \label{sp:integer}
\end{equation}
Indeed, the commutant on each multiplicity space is a real, complex or quaternionic matrix algebra. Its self-adjoint operator diagonalizes, and the variational principle selects the largest $k_\tau$ eigenvalues for a fixed multiplicity budget. Optimizing the resulting integer real-dimension budget proves \cref{sp:integer}. A nontrivial real $C_7$ rotation plane cannot be split into two separately admissible one-dimensional directions.

\begin{spectrallemma}[Unconstrained regression risk]
\label{sp:regression}
The optimal unconstrained matrices are $A_g=J^*U_gJ$, and
\begin{equation}
 \eta=\frac1{rm}\sum_g\norm{(\Id-P)U_gJ}_{\mathrm{HS}}^2
 =1-\operatorname{tr}(T^2)/r.
 \label{sp:eta}
\end{equation}
\end{spectrallemma}
\begin{proof}
Since $J$ is an isometry, $JA_g=PU_gJ$ is the least-squares projection of $U_gJ$ onto $\operatorname{range}J$. Averaging its residual squared gives $r-\operatorname{tr}(PT)$. Group averaging is the Hilbert--Schmidt orthogonal projection onto the commutant, so $\operatorname{tr}(PT)=\operatorname{tr}(T^2)$. Division by $r$ proves the result.
\end{proof}

\subsection{The fixed-dimensional risk bound}
\begin{spectraltheorem}[Risk of the entire exact orthogonal action class]
\label{sp:main}
Under the assumptions of \cref{sp:setup},
\begin{equation}
 \boxed{\quad 2(1-\gamma_r)\ \le\ L^\star\ \le\ 2(1-\gamma_r^2).\quad}
 \label{sp:sandwich}
\end{equation}
The lower bound holds for every $\rho\in\mathcal R_r(G)$, and some member of this same $r$-dimensional class satisfies the upper bound.
\end{spectraltheorem}
\begin{proof}
For a fixed $\rho$, define on maps $A:\R^r\to\mathcal H$
\[
 \Pi_\rho(A)=\frac1m\sum_gU_gA\rho_g^*.
\]
Each summand is a Hilbert--Schmidt isometry. Reindexing by inverses and products proves that the average is self-adjoint and idempotent; it projects onto the intertwiner space $\{A:U_gA=A\rho_g\ \forall g\}$. With $B=\Pi_\rho(J)$, expansion gives
\begin{equation}
 L(\rho)=2-\frac2r\langle J,B\rangle_{\mathrm{HS}}
        =2-\frac2r\norm B_{\mathrm{HS}}^2.
 \label{sp:projection-risk}
\end{equation}
Let $Q_B$ project onto $\operatorname{range}B$. This range is invariant, contained in $\mathcal W$, and has rank at most $r$. Cauchy--Schwarz yields
\[
 \norm B_{\mathrm{HS}}^2=\langle Q_BJ,B\rangle_{\mathrm{HS}}
 \le\norm{Q_BJ}_{\mathrm{HS}}\norm B_{\mathrm{HS}},\qquad
 \norm B_{\mathrm{HS}}^2\le\operatorname{tr}(PQ_B)
 =\operatorname{tr}(TQ_B)\le C_r(T).
\]
The zero-$B$ case is immediate. Substitution into \cref{sp:projection-risk} proves the lower bound.

For the upper bound, choose a maximizer $Q$ in \cref{sp:budget}, of rank $k\le r$. In an auxiliary space put
\[
 \widetilde{\mathcal H}=\mathcal H\oplus\R^{r-k},\quad
 \widetilde U_g=U_g\oplus\Id_{r-k},\quad
 \widetilde J=(J,0),\quad \widetilde Q=Q\oplus\Id_{r-k}.
\]
Let $V:\R^r\to\widetilde{\mathcal H}$ be an orthonormal basis of $\operatorname{range}\widetilde Q$. An SVD of $V^*\widetilde J$ and orthogonal Procrustes alignment give another such basis $J'$ with
\[
 \langle J',\widetilde J\rangle_{\mathrm{HS}}
 =\norm{V^*\widetilde J}_*
 \ge\norm{V^*\widetilde J}_F^2
 =\operatorname{tr}(\widetilde Q\widetilde J\widetilde J^*)=C_r(T).
\]
The inequality uses singular values in $[0,1]$, since $V$ and $\widetilde J$ are isometries. Invariance of $\operatorname{range}J'$ makes $\rho_g=J'^*\widetilde U_gJ'$ an exact orthogonal action on $\R^r$ and makes $J'$ an intertwiner. Hence, for $B'=\Pi_\rho(\widetilde J)$,
\[
 C_r(T)\le\langle J',\widetilde J\rangle_{\mathrm{HS}}
 =\langle J',B'\rangle_{\mathrm{HS}}
 \le\sqrt r\norm{B'}_{\mathrm{HS}}.
\]
Equation~\eqref{sp:projection-risk} gives $L(\rho)\le2-2C_r(T)^2/r^2$. The data have zero auxiliary coordinates, so this construction changes neither the original loss nor the predictor dimension; the final matrices remain $r\times r$.
\end{proof}

\paragraph{Relation to existing theory.}
\label{sp:positioning}
The averaging, intertwiner and dimension-accounting tools are standard in approximate representation theory \citep{gowers2017inverse,moore2015approximate}. The statement here connects invariant approximation to \emph{fixed-dimensional prediction risk}; it does not claim priority for those tools. Sun's realizability and error decomposition concern compatible transformations through an encoder \citep{sun2026transformationlaws}; EbC identifies jointly learned latent coordinates and actions under generative and diversity assumptions \citep{schmidt2025equivariance}. Neither observation contract is replaced by this frozen-coordinate bound.

\subsection{Finite observable form and an insufficient-summary example}
\label{sp:observable}
Let $D=mr$, $v(c,s)=m^{-1/2}(x_{c,g\cdot s})_{g\in G}\in\R^{1\times D}$ and $K=\mathbb E[v^\top v]$. Its blocks are $K_{g,h}=m^{-1}\mathbb E[x_g^\top x_h]$. For $a=(a_g)_g$, define
\[
 Wa=m^{-1/2}\sum_gU_gJa_g,\qquad (\Lambda_ta)_k=a_{kt^{-1}}.
\]
Direct multiplication gives $WW^*=T$, $W^*W=K$, and $U_tW=W\Lambda_t$. Thus $K$ commutes with the known block action $\Lambda$, has trace $r$, and is a positive contraction. The polar decomposition of $W$ identifies $\operatorname{range}K$ with $\mathcal W$ equivariantly: their nonzero spectra agree \emph{with their real irreducible types}. Therefore the optimization in \cref{sp:budget}, now using $K,\Lambda$ on $\R^D$, gives
\[
 C_r(K)=C_r(T),\qquad \eta=1-\operatorname{tr}(K^2)/r.
\]
Zero-eigenvalue blocks add no trace objective. This construction uses complete same-context labeled orbits as a structural diagnostic; its information is not the input available to a single-activation predictor. Computing the statistic still requires the type-constrained optimization, not ordinary top-$r$ eigenvalue truncation.

\paragraph{Equal regression risk and eigenvalues can conceal different action risk.}
\label{sp:counterexample}
Take $G=S=C_2\times C_2\times C_7$ acting by addition, uniform states $(a,b,k)$, one context, and $r=1$. Set
\[
 x_A(a,b,k)=\frac{(-1)^a+(-1)^b}{\sqrt2},\qquad
 x_B(a,b,k)=\sqrt2\cos(2\pi k/7).
\]
Both have zero mean and unit second moment. For $g=(u,v,w)$ their optimal scalar regression coefficients are
\[
 a_A(g)=\tfrac12\big((-1)^u+(-1)^v\big),\qquad
 a_B(g)=\cos(2\pi w/7).
\]
Both satisfy $\mathbb E_g a(g)^2=1/2$, hence $\eta=1/2$. For $A$, $T$ has eigenvalue $1/2$ on each of two distinct real one-dimensional characters; for $B$, it is $\Id_2/2$ on a real irreducible rotation plane. The nonzero eigenvalues are thus the same, $\{1/2,1/2\}$. An exact real one-dimensional action is a character $\chi:G\to\{\pm1\}$ and is necessarily trivial on $C_7$. Its loss is $2-2\mathbb E_g[\chi(g)a(g)]$. Character orthogonality makes the largest correlation $1/2$ for $A$ and zero for $B$, so $L_A^\star=1$ and $L_B^\star=2$; correspondingly $\gamma_1=1/2$ and $0$. This illustrates the standard representation-dimension obstruction: regression risk and an untyped spectrum are insufficient summaries. The complete labeled observations distinguish the examples, so this is not observational nonidentifiability of the full data.

\paragraph{Scale of the type-budget correction for C7 and D4.}
\label{sp:type-scale}
For these two groups all real irreducible types have dimension one or two. Define the ordinary spectral budget $C_r^{\mathrm{untyped}}(K)$ as the sum of the largest $r$ eigenvalues of $K$, including zeros if needed. Order tied eigenvectors by complete irreducible blocks. This top-$r$ selection can split at most one two-dimensional block at its boundary. Dropping the selected half yields a feasible invariant subspace of dimension at most $r$ and loses at most one eigenvalue, which is at most one because $K\preceq\Id$. Hence
\[
 0\le \Delta_{\mathrm{type}}
 :=C_r^{\mathrm{untyped}}(K)-C_r(K)\le1,
 \qquad \Delta_{\mathrm{type}}/r\le1/r.
\]
This elementary consequence gives a normalized upper bound of $1/16$ at $r=16$; it is not an observed effect or a claim that the bound is tight. A zero type correction would leave the class-risk bound valid. The one-dimensional separation above therefore does not predict a large correction on the LLM panel, and does not justify searching dimensions for a positive correction.

For the four-corner action $R(i)=i+1$, $S(i)=1-i\pmod4$, the state permutation representation decomposes as $\chi_{+,+}\oplus\chi_{-,-}\oplus E$, where $\chi_{a,b}(R)=a$, $\chi_{a,b}(S)=b$ and $E$ is the real two-dimensional type. Indeed, the stabilizer of corner zero is $\{e,R^3S\}$, so a one-dimensional type has multiplicity $(1+ab)/2$. Context copies do not introduce the absent types: $K$ has zero energy in $\chi_{+,-}$ and $\chi_{-,+}$. This is an analytical input check for any future numerical diagnostic, not a completed numerical check. The corner action has four states even though the group has eight elements.

\subsection{A finite-context certificate and its practical limitation}
\label{sp:certificate}
Suppose $C_1,\ldots,C_n$ are independent identically distributed complete context blocks, independent of the choice of coordinates and preprocessing. Assume the \emph{population} bound $\norm{z_{C,s}}_2\le H$ almost surely for every state, with $H$ fixed in advance. Define
\[
 \Sigma_i=p^{-1}\sum_sz_{C_i,s}^\top z_{C_i,s},\quad
 v_z(c,s)=m^{-1/2}(z_{c,g\cdot s})_g,\quad
 M_i=p^{-1}\sum_sv_z(C_i,s)^\top v_z(C_i,s),
\]
and sample means $\widehat\Sigma,\widehat M$ with population means $\Sigma,M$. When $\widehat\Sigma\succ0$, let
\[
 \widehat K=(\Id_m\otimes\widehat\Sigma^{-1/2})\widehat M
                    (\Id_m\otimes\widehat\Sigma^{-1/2}).
\]
Complete uniform state averaging preserves the fixed $\Lambda$ symmetry exactly and gives $\operatorname{tr}\widehat K=r$.

\begin{spectraltheorem}[Conservative complete-block certificate]
\label{sp:finite}
For $\alpha\in(0,1)$, set
\[
 \tau_n=H^2\sqrt{\frac8n\log\frac{2(D+r)}\alpha},\qquad
 a_n=\lambda_{\min}(\widehat\Sigma)-\tau_n.
\]
If $a_n>0$, define $\varepsilon_n=\tau_n/a_n+H^2\tau_n/a_n^2$, $\widehat\gamma_r=C_r(\widehat K)/r$, and
\[
 \gamma_- =\max\{0,\widehat\gamma_r-\varepsilon_n\},\qquad
 \gamma_+ =\min\{1,\widehat\gamma_r+\varepsilon_n\}.
\]
There is an event of probability at least $1-\alpha$ on which, whenever this precheck passes,
\begin{equation}
 2(1-\gamma_+)\le L^\star\le2(1-\gamma_-^2).
 \label{sp:confidence}
\end{equation}
If $a_n\le0$, the procedure returns inconclusive; discarding weak directions changes the specified $r$-dimensional target. This bounds the unconditional probability of issuing an incorrect interval, not the error probability conditional on passing the precheck.
\end{spectraltheorem}
\begin{proof}
Each $\Sigma_i,M_i$ lies between zero and $H^2\Id$, so its centered square is bounded by $H^4\Id$. Applying the two-sided matrix Hoeffding inequality \citep[Theorem~1.3]{tropp2012user} and a union bound for dimensions $r,D$ gives, with probability at least $1-\alpha$,
\[
 \norm{\widehat\Sigma-\Sigma}_{\rm op}\le\tau_n,\qquad
 \norm{\widehat M-M}_{\rm op}\le\tau_n.
\]
On this event and when $a_n>0$, both second moments are at least $a_n\Id$. For positive $A,B\succeq a\Id$, the inverse-square-root resolvent formula gives
\[
 \norm{A^{-1/2}-B^{-1/2}}_{\rm op}
 \le\frac{\norm{A-B}_{\rm op}}\pi\int_0^\infty\frac{t^{-1/2}}{(a+t)^2}\,dt
 =\frac{\norm{A-B}_{\rm op}}{2a^{3/2}}.
\]
Writing $A_0=\Id_m\otimes\Sigma^{-1/2}$ and $\widehat A_0=\Id_m\otimes\widehat\Sigma^{-1/2}$, expand
\[
 K-\widehat K=A_0(M-\widehat M)A_0
 +(A_0-\widehat A_0)\widehat M A_0
 +\widehat A_0\widehat M(A_0-\widehat A_0).
\]
Using $\norm{\widehat M}_{\rm op}\le H^2$ yields $\norm{K-\widehat K}_{\rm op}\le\varepsilon_n$. Every feasible invariant projection has rank at most $r$, so
\[
 |\operatorname{tr}((K-\widehat K)Q)|\le r\varepsilon_n,
 \qquad |C_r(K)-C_r(\widehat K)|\le r\varepsilon_n.
\]
The feasible set is fixed by $G,\Lambda,r$, not selected from the data. Apply \cref{sp:main} to obtain \cref{sp:confidence}.
\end{proof}

\paragraph{Why this does not certify the present experiments.}
The bound is formally valid but can be uninformative. Necessarily
\[
 \lambda_{\min}(\widehat\Sigma)\le\operatorname{tr}(\widehat\Sigma)/r\le H^2/r,
 \qquad a_n>0\ \Longrightarrow\
 n>8r^2\log\frac{2r(m+1)}\alpha.
\]
At $r=16$ and $\alpha=0.05$, even the possibility of passing the precheck requires at least 17,492 independent complete contexts for $C_7$, or 17,734 for $D_4$. These are necessary conditions for this conservative formula, not sufficient conditions, intrinsic sample-complexity lower bounds or proposed experimental sizes. Even after passing, \cref{sp:confidence} may be $[0,2]$.

The current 40-template panel and its overlapping partitions cannot meet this condition. States, actions, layers and repeated splits are not additional independent contexts. A sample maximum cannot replace the stipulated population bound $H$; clipping would change the target. Data-dependent preprocessing requires independent fitting data or a separate analysis. Multiple reported model/layer targets require an explicit simultaneous-error allocation; selecting the best layer is not covered by the single-target formula. The complete observed panel could instead define a finite empirical population and an exact descriptive version of \cref{sp:sandwich}, but that would not establish new-context coverage. No such statistic or certificate is reported as an experimental result here. The retained scores, failed gates and inferential scope are unchanged.

\input{sections/methods.tex}
\input{sections/results.tex}
\input{sections/characterization.tex}
\input{sections/reproduction.tex}
\input{sections/prospective.tex}
\input{sections/related_work.tex}
\end{document}
